% W5i73_HardProblems.tex
% DFD 20-Agent Research Programme --- Iteration 73 (i73)
% Wave 5 Cycle (i52--i100): TWENTY-SECOND ITERATION
% Hard Problems Register
% Date: 2026-03-24

\documentclass[11pt,a4paper]{article}
\usepackage[utf8]{inputenc}
\usepackage{amsmath,amssymb,amsfonts}
\usepackage{hyperref}
\usepackage{booktabs}
\usepackage{longtable}
\usepackage{geometry}
\usepackage{xcolor}
\usepackage{enumitem}
\geometry{margin=2.5cm}

\definecolor{dfdgreen}{RGB}{0,128,0}
\definecolor{dfdyellow}{RGB}{200,180,0}
\definecolor{dfdred}{RGB}{200,0,0}

\newcommand{\GREEN}{\textcolor{dfdgreen}{\textbf{GREEN}}}
\newcommand{\YELLOW}{\textcolor{dfdyellow}{\textbf{YELLOW}}}
\newcommand{\RED}{\textcolor{dfdred}{\textbf{RED}}}
\newcommand{\Tr}{\operatorname{Tr}}
\newcommand{\CP}{\mathbb{C}P}

\title{\Large W5i73: Hard Problems Register\\[6pt]
\normalsize Wave 5 Cycle: Twenty-Second Iteration (i73 of i52--i100)\\[6pt]
3 YELLOW $\cdot$ Integration Residuals $\cdot$ N-Body Convergence $\cdot$ $E_G$ Systematics}
\author{DFD 20-Agent Research Programme}
\date{2026-03-24}

\begin{document}
\maketitle

%=============================================================================
\section{Hard Problem Status Update}
%=============================================================================

Three YELLOW items persist from i72. The v3.3 integration pass (HP-5 from i72) is now resolved. New hard problems emerge from N-body and $E_G$ work.

%=============================================================================
\section{HP-1: $\alpha$ Residual ($0.56$ ppm) --- YELLOW (Unchanged)}
%=============================================================================

\subsection{Status Update}
No change from i72. The residual remains at $+0.763 \times 10^{-6}$ ($0.56$ ppm). The integration pass confirmed $\alpha^{-1} = 137.035\,999\,847$ throughout all 20 v3.3 sections. Error budget unchanged: Toeplitz truncation dominates at $0.42$ ppm.

\subsection{New Development: $O(k^{-2})$ Correction Estimate}
Agent 17 computed the next-order truncation correction:
\[
\delta\alpha^{-1}_{k^{-2}} = \frac{\pi^{3/2}}{24}\Tr(Y^2)\cdot\frac{1}{8k^2}\left(1 + \frac{3}{2k}\right) = \frac{\pi^{3/2}\cdot 10}{24}\cdot\frac{1}{8\cdot 3600}\cdot 1.025 = 0.0059
\]
This gives $\delta(\alpha^{-1})/\alpha^{-1} = 0.043$ ppm. Including this:
\[
\alpha^{-1}(k=60,\, O(k^{-2})) = 137.035\,999\,847 - 0.0059 = 137.035\,999\,841
\]
Residual from CODATA: $+0.757 \times 10^{-6}$ ($0.55$ ppm). Marginal improvement ($0.56 \to 0.55$ ppm). The $O(k^{-2})$ correction is subdominant. \textbf{Full resolution requires electroweak threshold matching or $k \to \infty$ continuum theory.}

%=============================================================================
\section{HP-2: $\Sigma m_\nu = 61.4$ meV --- YELLOW (Unchanged)}
%=============================================================================

No change from i72. Phase 0B self-consistency correction is negligible ($+0.02$ meV). Margin from cosmological bound: 2.2 meV. Strategy: await JUNO ($\sim 2031$) and Euclid + CMB-S4 ($\sigma(\Sigma m_\nu) \sim 15$ meV by 2028).

%=============================================================================
\section{HP-3: $E_G$ Prediction --- YELLOW (Advancing)}
%=============================================================================

\subsection{Status Update}
$E_G$ paper now at 70\%. Euclid forecast section complete. Key advance: systematic error analysis quantified.

\subsection{Systematic Error Budget for $E_G$}
\begin{center}
\begin{tabular}{lcc}
\toprule
\textbf{Systematic} & \textbf{Fractional effect on $\epsilon$} & \textbf{Mitigation} \\
\midrule
Photo-$z$ errors & 15\% degradation & Spectroscopic calibration \\
Intrinsic alignments & $< 5\%$ at $z > 0.5$ & NLA model subtraction \\
Baryonic feedback & $< 3\%$ at $k < 0.3\,h$/Mpc & Conservative $k$-cut \\
Magnification bias & $< 2\%$ & Slope measurement \\
Source clustering & $< 1\%$ & Random subtraction \\
\midrule
\textbf{Total systematic floor} & $\sim 16\%$ & \\
\bottomrule
\end{tabular}
\end{center}

With systematics, the DFD signal significance degrades from $3.1\sigma$ (stat only) to $2.6\sigma$ (stat + sys). This is still a meaningful detection. \textbf{Recommendation}: quote $2.6\sigma$ as the realistic detection significance.

\subsection{New Finding: Redshift-Dependent Systematic Cancellation}
The DFD signature $\epsilon(z)$ is monotonically rising, while most systematics (photo-$z$, IA) are redshift-independent or declining. This means the highest-$z$ bins ($z > 1.0$) have the best signal-to-systematic ratio. Agent 10's Fisher analysis confirms: $z = 1.3$ bin alone gives $1.8\sigma$, while $z = 0.5$ bin gives only $0.9\sigma$. \textbf{High-$z$ bins are key.}

%=============================================================================
\section{HP-4: N-Body Convergence (NEW)}
%=============================================================================

\subsection{Problem Statement}
The $512^3$ validation run shows $P^{\rm DFD}/P^{\rm \Lambda CDM} = 1.008 \pm 0.003$ at $k = 0.3\,h$/Mpc. But the error bar is dominated by cosmic variance at the box size $L = 250\,h^{-1}$Mpc. The full $1024^3$ run ($L = 500\,h^{-1}$Mpc) should reduce this by $\sim\sqrt{8}$ (factor of 2.8). Expected: $\sigma(P^{\rm DFD}/P^{\rm \Lambda CDM}) \sim 0.001$ at $k = 0.3\,h$/Mpc.

\subsection{Convergence Requirements}
\begin{enumerate}[nosep]
\item Mass resolution: $m_p = 8.5 \times 10^9\,h^{-1} M_\odot$ at $1024^3$ in $L = 500\,h^{-1}$Mpc. Sufficient for halo mass function down to $M > 10^{12}\,h^{-1} M_\odot$.
\item Force resolution: softening $\epsilon = 10\,h^{-1}$kpc. Sufficient for $k < 3\,h$/Mpc.
\item Time integration: adaptive timestep with Courant factor 0.3. Energy conservation: $|\Delta E/E| < 10^{-4}$ over full run.
\item $\rho$-field solver: PM grid with $N_{\rm grid} = 2048^3$ (twice particle resolution). Convergence test: $1024^3$ grid gives $P(k)$ within 0.2\% of $2048^3$ at $k < 1\,h$/Mpc.
\end{enumerate}

\subsection{Timeline}
Full run: 2 weeks. Post-processing: 1 week. Analysis: 2 weeks. First results expected by $\sim$i78.

%=============================================================================
\section{HP-5: v3.3 Integration Residuals}
%=============================================================================

\subsection{Problem Statement}
Despite the integration pass, some residual issues remain for submission-readiness:

\begin{enumerate}[nosep]
\item \textbf{Page count}: 218 pages + 22 appendix pages = 240 total. Physics Reports allows up to 300 pages, so acceptable, but referee burden is high. Consider: supplemental material separation or two-part submission.
\item \textbf{Self-citation ratio}: v3.3 cites 287 references, of which 23 are DFD-authored (8.0\%). This is within acceptable bounds (threshold concern at $> 15\%$), but some sections (17, 19) have higher local ratios. Mitigated by adding external references in those sections.
\item \textbf{Code availability}: v3.3 references ``DFD software package'' (JOSS submission-ready) but JOSS paper not yet submitted. For submission, need either JOSS DOI or Zenodo DOI for code.
\end{enumerate}

\subsection{Resolution}
Items 1--2 are minor and will be addressed in i74 (submission-ready pass). Item 3 requires coordination with JOSS submission timeline. Recommendation: submit v3.3 to Physics Reports with Zenodo DOI; submit JOSS simultaneously.

%=============================================================================
\section{HP-6: DESI Y3 Preparation}
%=============================================================================

\subsection{Problem Statement}
DESI Y3 data expected late 2026 or Q1 2027. DFD paper (now 100\%) predicts $w_a = -0.12 \pm 0.05$, distinguishable from $\Lambda$CDM at $3.2\sigma$ with Y3. Need to prepare analysis pipeline for rapid confrontation.

\subsection{Required Preparation}
\begin{enumerate}[nosep]
\item Pre-compute DFD $w_0$--$w_a$ likelihood surface on fine grid ($100 \times 100$).
\item Prepare MCMC chains with DFD theory vector.
\item Draft ``DFD confrontation with DESI Y3'' paper template (results-ready, just awaiting data).
\end{enumerate}

\subsection{Timeline}: i74--i76.

%=============================================================================
\section{Hard Problems Summary Table}
%=============================================================================

\begin{center}
\begin{tabular}{clccc}
\toprule
\textbf{\#} & \textbf{Problem} & \textbf{Status} & \textbf{Priority} & \textbf{Timeline} \\
\midrule
HP-1 & $\alpha$ residual 0.55 ppm & \YELLOW & High & Ongoing \\
HP-2 & $\Sigma m_\nu$ margin & \YELLOW & High & 2031 (JUNO) \\
HP-3 & $E_G$ prediction & \YELLOW & High & Oct 2026 \\
HP-4 & N-body convergence & \GREEN & High & $\sim$i78 \\
HP-5 & v3.3 residuals & \GREEN & Immediate & i74 \\
HP-6 & DESI Y3 preparation & \GREEN & Medium & i74--i76 \\
\bottomrule
\end{tabular}
\end{center}

\end{document}
